\chapter{Memory access and locality}
\MARSChapterHeading{Memory access and locality}
\label{revision:locality}

Two loops can copy the same words to the same final positions while visiting
addresses in different orders. Those orders determine which nearby words are
used together and how long a cache line must remain available before it is used
again. Correct output is the starting requirement; explaining locality requires
an address trace as well.

The complete \texttt{matrix-locality.asm} program in \texttt{revision/labs/files}
offers two ways to transpose a matrix: an ordinary row traversal and
a traversal grouped into four-by-four tiles. Each launch selects one method and
runs it on a fresh input/output pair. Both implement original project
\texttt{011-matrix-transpose}'s contract. Neither changes the word values.

This chapter assumes indexed arrays, nested loops, full-width pointers and
function preservation. You should be able to derive the source and destination
addresses, explain a line hit or miss in the stated model, and compare measurements
without treating them as elapsed-time predictions. The independent exercise rotates
a rectangular matrix clockwise, with the same separation between data mapping
and the order chosen to perform it.

\section{Derive two row widths before writing a loop}

A row-major matrix with $R$ rows and $C$ columns stores source cell $(i,j)$ at
linear index $iC+j$. Transpose places that value at output coordinate $(j,i)$.
The output has $C$ rows and $R$ columns, so its linear index is $jR+i$.
For four-byte words the two addresses are therefore
$input+4(iC+j)$ and $output+4(jR+i)$. The destination row width is $R$, not $C$.

Use the default three-by-five source, filled with values zero through fourteen.
Source $(1,0)$ contains five at source index five. It goes to output $(0,1)$,
index one. Source $(0,1)$ contains one and goes to output $(1,0)$, index three.
These two examples separate swapping coordinates from merely copying the same
linear index. A square-only test can hide a mistaken destination row width.

\begin{figure}[htbp]
  \centering
  \begin{tikzpicture}[x=7mm,y=-7mm,alt={A three by five source contains rows zero through four, five through nine, and ten through fourteen. Its five by three transpose contains rows zero five ten; one six eleven; two seven twelve; three eight thirteen; four nine fourteen. Shading covers source columns zero through three in all three rows and the corresponding first four complete destination rows. This is the first clipped four by four source tile, with twelve cells.}]
    \foreach \i in {0,1,2} {\foreach \j in {0,1,2,3} {\fill[black!12] (\j,\i) rectangle ++(1,1);}}
    \foreach \i in {0,1,2,3} {\fill[black!12] (7,\i) rectangle ++(3,1);}
    \foreach \x in {0,...,5} {\draw (\x,0)--(\x,3);}
    \foreach \y in {0,...,3} {\draw (0,\y)--(5,\y);}
    \foreach \x in {7,...,10} {\draw (\x,0)--(\x,5);}
    \foreach \y in {0,...,5} {\draw (7,\y)--(10,\y);}
    \foreach \i in {0,1,2} {\foreach \j in {0,...,4} {\pgfmathtruncatemacro{\n}{5*\i+\j}\node[font=\small] at (\j+.5,\i+.5) {\n};}}
    \foreach \i in {0,...,4} {\foreach \j in {0,1,2} {\pgfmathtruncatemacro{\n}{5*\j+\i}\node[font=\small] at (7+\j+.5,\i+.5) {\n};}}
    \node[anchor=south,font=\small\bfseries] at (2.5,-.2) {Source: $3\times5$};
    \node[anchor=south,font=\small\bfseries] at (8.5,-.2) {Output: $5\times3$};
    \draw[->,thick] (5.35,1.4)--(6.65,1.4);
  \end{tikzpicture}
  \caption{Transpose changes row width as well as coordinates. The source values name their original linear indices. Shading identifies the first clipped tile; the remaining source column becomes the final output row.}
  \label{fig:revision-locality-mapping}
\end{figure}


Each dimension may be zero through sixteen. The largest active matrix has 256
words and the largest valid relative displacement is 1020 bytes. These bounds
make word arithmetic sufficient for index products, but do not make a base address
32 bits wide. Add displacements to pointers with full-width address arithmetic.
A one-kilobyte allocation can still live above the first four gigabytes of memory.

The source and destination are separate, nonoverlapping active windows inside
256-word allocations. Passed pointers can name interior windows; do not replace
them with global label addresses. Protect source words, guards and inactive output
prefixes/suffixes from stores. An empty dimension means no matrix reads or writes.
A zero return count alone would not prove that the function avoided an invalid
first load before noticing the empty shape.

The values are arbitrary 32-bit patterns. A \texttt{lw} can sign extend the loaded
word into a register, but \texttt{sw} stores its same low 32 bits. Since transpose
never performs arithmetic on that value, signed interpretation does not alter the
copied pattern. A byte or halfword store would fail this requirement.

\section{Follow the ordinary address recurrence}

The ordinary function visits source rows in increasing order and each row's
columns in increasing order. Within one row, advancing $j$ by one advances the
source pointer by four bytes. In the destination, that same advance changes
$jR+i$ by $R$, so its pointer advances by $4R$ bytes. After a complete source row,
the destination's next starting position advances by one word.

The complete function checks both empty cases and returns $RC$ in
\texttt{\$v0}. It prepares destination stride $4R$ in \texttt{\$v1}, row index
zero in \texttt{\$t0}, the source cursor in \texttt{\$t2}, and destination row
origin in \texttt{\$t3}. Its loop is:

\noindent\begin{minipage}{\linewidth}
\MARSCodeLabel{Read a source row while writing a destination column}
\begin{minted}{mars-idea}
naive_row:
    move $t1, $zero
    move $t4, $t3
naive_cell:
    lw $t5, 0($t2)
    sw $t5, 0($t4)
    daddiu $t2, $t2, 4
    daddu $t4, $t4, $v1
    addiu $t1, $t1, 1
    bltc $t1, $a3, naive_cell
    nop
    daddiu $t3, $t3, 4
    addiu $t0, $t0, 1
    bltc $t0, $a2, naive_row
    nop
\end{minted}
\end{minipage}

The source cursor never resets: the next source row follows the preceding one
contiguously. The destination cursor \texttt{\$t4} is reset from the next row
origin, because its inner-loop jumps do not lead to that origin. Keeping these
two roles separate avoids accumulating the destination stride into the wrong row.
The function is a leaf and leaves saved registers and the stack pointer untouched.

For $R=3$, the first five source byte displacements are $0,4,8,12,16$ and the
corresponding destination displacements are $0,12,24,36,48$. Values one through
four therefore do not fill the next four consecutive output words. Derive the
next source row's destination displacements $4,16,28,40,52$ before running Debug.
All fifteen source words must appear once in the five-by-three output.

\section{Explain the cache model one access at a time}

The Data Cache Simulator groups addresses into lines, called blocks in its
controls. Here each line contains four words, or sixteen bytes, and the cache
has eight lines: 128 bytes total. A memory line number is $\lfloor address/16\rfloor$.
An access is a hit if that line is already recorded in an allowed cache entry;
otherwise it is a miss and installs that line, possibly replacing another.
The initial cold cache has no resident lines.

Nearby words can share a line. Using them close together is spatial locality.
Returning to a word or resident line before it is replaced provides reuse over
time. In this transpose each individual word is read or written only once;
the useful reuse is between different words in the same line. The model records
line presence, not cached data values. Its architectural memory operations still
produce the same results regardless of hits and misses.

In fully associative placement, any of the eight entries can hold any line.
The least-recently-used policy, LRU, replaces the entry whose last access occurred
furthest in the past when all allowed entries are occupied. A hit updates that
recency too. This model allocates a line on a store miss as well as a load miss.
An output miss is therefore part of the experiment, not a reason to ignore stores.

\begin{table}[htbp]
\centering
\small
\begin{tabular}{rrrrll}
\toprule
Source $(i,j)$ & Input byte offset & Output byte offset & Value & Read & Write\\
\midrule
$(0,0)$ & 0 & 0 & 0 & miss & miss\\
$(0,1)$ & 4 & 12 & 1 & hit & hit\\
$(0,2)$ & 8 & 24 & 2 & hit & miss\\
$(0,3)$ & 12 & 36 & 3 & hit & miss\\
$(0,4)$ & 16 & 48 & 4 & miss & miss\\
\bottomrule
\end{tabular}
\caption{The first source row of the ordinary $3\times5$ transpose in a cold, fully associative eight-line cache with sixteen-byte lines. Offsets are relative to their separate aligned input and output bases. Each row executes its read before its write.}
\label{tab:revision-locality-first-row}
\end{table}


At the first input read, source bytes zero through fifteen belong to the first
line, so the later reads at offsets four, eight and twelve hit. The output's first
store installs its own first line; offset twelve still belongs to that line and
hits. Output offset 24 belongs to a different line and misses. Source and output
are separate allocations even when their relative offsets match. Across the
complete three-by-five case, each occupies four lines; all eight fit, yielding
eight initial-touch misses and no further misses in this fully associative model.

The tool observes forward simulated data-access notices. These aligned word loads
and stores each contribute one notice. Instruction fetches and external debugger
inspection do not contribute. Recording a line does not mean the program executed
four separate loads, nor does the tool calculate fill latency or byte traffic.
It does not model a processor pipeline, prefetching, write-back traffic or host
execution time. Architectural backstep keeps cache observations rather than
rewinding their history. Start a fresh run, or reset at the chosen stopped entry,
to establish a new cold comparison.

\section{Complete a small tile before moving elsewhere}

The ordinary destination jumps can separate uses of nearby destination words by
many intervening accesses. Tiling changes that schedule. Choose a four-by-four
source tile, finish its source reads and corresponding output writes, then move
to the next tile. The mathematical mapping remains $(i,j)\mapsto(j,i)$.

For tile origin $(b_i,b_j)$, the exclusive ends are
$e_i=\min(b_i+4,R)$ and $e_j=\min(b_j+4,C)$. Iterate $i$ from $b_i$ to just before
$e_i$ and $j$ from $b_j$ to just before $e_j$. The last tiles can be narrower or
shorter than four; their clipped bounds are part of correctness, not a cleanup
optimization to add after testing only square multiples of four.

The tiled routine holds the tile origin in \texttt{\$t0}/\texttt{\$t1} and
computes its ends as follows. The rest of the complete function prepares row
pointers from $iC+b_j$ and $b_jR+i$, then uses the same four-byte source and
$4R$-byte destination increments inside that row.

\noindent\begin{minipage}{\linewidth}
\MARSCodeLabel{Clip both exclusive tile ends to the actual shape}
\begin{minted}{mars-idea}
tile_begin:
    addiu $t2, $t0, 4
    blec $t2, $a2, tile_rows_ready
    nop
    move $t2, $a2
tile_rows_ready:
    addiu $t3, $t1, 4
    blec $t3, $a3, tile_columns_ready
    nop
    move $t3, $a3
\end{minted}
\end{minipage}

In the three-by-five example, the first tile contains three rows and four columns,
or twelve cells. Its output covers linear indices zero through eleven, as shaded
in Figure~\ref{fig:revision-locality-mapping}. The final tile contains the remaining
three source cells and fills output indices twelve through fourteen. Its source
IDs are four, nine and fourteen: the tile is not a contiguous twelve-word source
prefix followed by its remainder.

For a sixteen-by-sixteen matrix, a four-by-four source tile touches four source
lines and four destination lines. In the fully associative eight-line model,
those eight lines can coexist. By completing their sixteen cells together, the
function uses each line's words before leaving that tile. The whole matrix then
needs one initial touch for each of its 64 source lines and 64 destination lines.
The observed 128 misses reach that lower bound under this allocation and policy.

Placement can prevent lines from coexisting despite enough total capacity.
Direct mapping has eight one-entry sets, selecting set $line\bmod8$. Two-way
placement has four two-entry sets, selecting $line\bmod4$. A line can occupy only
its selected set; LRU chooses within that set, not across the whole cache.
Its tag identifies which of the different lines mapping to that set is present.

The public output starts 1088 bytes, or 68 lines, after input. For sixteen rows,
consecutive destination rows are 64 bytes, or four lines, apart. Their direct-map
sets alternate between two numbers; in the two-way model they all select the
same set. A tile's eight live lines therefore cannot necessarily occupy eight
different entries. Unused entries in other sets do not help a crowded set.
Derive these set numbers from the actual bases before attributing the difference
to the word values or the total matrix size.

\section{Measure the scope and the competing costs}

The public measurement tool creates fresh sessions and cold cache observations
at function entry. It stops at \texttt{transpose\_done}, before the caller stores
the returned count. It records completed simulator instructions and data notices
for the whole function. A stack-using implementation would include its auxiliary
traffic in that scope; these two leaf implementations have none.

\begin{table}[htbp]
  \centering
  \small
  \setlength{\tabcolsep}{8pt}
  \begin{tabular}{llrrrrr}
    \toprule
    Shape        & Method   & Instructions & Accesses & Full misses & Direct & Two-way \\
    \midrule
    $3\times5$   & Ordinary & 119          & 30       & 8           & 8      & 8       \\
                 & Tiled    & 197          & 30       & 8           & 8      & 8       \\
    $8\times8$   & Ordinary & 443          & 128      & 80          & 86     & 80      \\
                 & Tiled    & 625          & 128      & 32          & 48     & 54      \\
    $16\times16$ & Ordinary & 1643         & 512      & 320         & 344    & 320     \\
                 & Tiled    & 2457         & 512      & 128         & 344    & 320     \\
    $15\times13$ & Ordinary & 1271         & 390      & 244         & 230    & 237     \\
                 & Tiled    & 2059         & 390      & 174         & 227    & 235     \\
    \bottomrule
  \end{tabular}
  \caption{Measurements of the exact public worked program. Every cache starts cold, has eight lines of four words and uses LRU within a set. The function interval ends before the caller's count store. Full, direct and two-way placements have the same 128-byte capacity.}
  \label{tab:revision-locality-measurements}
\end{table}


Both methods perform two accesses per cell, so sixteen-by-sixteen gives 512
notices. Fully associative misses fall from 320 ordinary to 128 tiled, but the
instruction count rises from 1,643 to 2,457. Tile bounds, repeated pointer formation
and extra loop levels cost instructions. Neither column alone establishes a
speedup: the tool supplies no miss penalty, cycle model or elapsed host time.
A different implementation or machine requires its own justified measurement.

Direct mapping gives 344 misses for both methods, and two-way placement gives
320 for both at this shape. Tiling improved the fully associative schedule but
did not remove these placement conflicts. Even the ordering between direct and
two-way counts need not be the same for every trace: at eight-by-eight, the tiled
counts are 48 and 54. Changing ways also changes the number of sets and the lines
that compete. Read the actual geometry and trace rather than treating a larger
way count as a universal promise for every LRU workload.

The smaller three-by-five matrices already fit the fully associative capacity,
so tiling adds instructions without reducing their eight misses. Fifteen-by-thirteen
has partially filled edge lines and clipped tiles; use its actual dimensions
rather than extending the square's 128-miss argument to it.

Repeat the experiment from the public repository root with JDK 25:
\begin{verbatim}
python3 tools/measure_revision_transpose.py \
  --java /path/to/jdk-25/bin/java \
  --cli /path/to/cli.jar --output /new/report/directory
\end{verbatim}
The tool checks the result and retains exact source/CLI hashes, measurements and
the fully associative access log. A new output directory keeps an older receipt
from being overwritten. Compare dimensions, base addresses, policy, initial cache
state and interval before comparing numbers. Assembly, session startup and host
time are outside these instruction/access counts.

\section{Derive a different mapping independently}

Implement \texttt{rotate\_clockwise} in \texttt{matrix-rotate.asm}. A clockwise
quarter-turn maps source $(i,j)$ to output $(j,R-1-i)$. The output still has
$C$ rows and $R$ columns, but one axis is reversed. Its linear destination index
is $jR+(R-1-i)$, rather than transpose's $jR+i$.

For source rows $[0,1,2]$ and $[3,4,5]$, the two-by-three matrix rotates into
three rows $[3,0]$, $[4,1]$, $[5,2]$. Follow two corners to establish direction
before deriving your loops. Ordinary transpose would produce $[0,3]$, $[1,4]$,
$[2,5]$ instead. A one-by-one test cannot distinguish those mappings.

Use the same full-pointer, nonoverlapping-window and zero-through-sixteen dimension
preconditions. Return $RC$; preserve saved registers, the stack pointer, caller
storage, source and guards. Each active source word receives one aligned read
and each corresponding destination one aligned write, with no destination reads
or inactive stores. Choose any traversal order; this independent task has no
four-by-four tile schedule requirement.

You might scan source rows and derive reversed destination origins, or scan output
rows and derive source addresses backward. Bound relative indices separately from
full addresses, and prove that every legal coordinate reaches exactly one output
cell. Then use empty, rectangular and interior-window tests. The written lab
supplies checks and investigations without prescribing your private labels or
register allocation.
